%----------------------------------------------------------------------
\section{Temporal Multiplexing: The XOR-Bridged Register}
\label{sec:time-sharing}\footnote{Section-level Toulmin. \emph{Move}: introduce the second compound production, which reuses a single register across $\GF(2)$-linear atoms by bridging them with XOR save/restore. \emph{[DAF] comparison}: [DAF]'s read/write path asymmetry distinguishes the construction of types ($\kappa$-evolution) from their elimination (evaluation); this section's XOR-bridge is the scheduling-level read/write asymmetry that lets the same register carry multiple atomic computations in succession. \emph{Residue instance}: the XOR bridge itself is a residue --- a small scratch value that records which lanes were not touched by the atom and restores them after.}
%----------------------------------------------------------------------

\subsection{The XOR-bridge}

Over $\GF(2)$, every operation is reversible because
$x \oplus m \oplus m = x$ for any mask $m$. The
\emph{XOR-bridge} exploits this reversibility to run
an atom $L$ over part of a register, then undo $L$'s
effect on the remaining lanes, leaving the register
free for the next atom.

\begin{definition}[XOR-bridge micro-kernel]
  \label{def:xor-bridge}
  For a $\GF(2)$-linear atom $L$ and a mask $m$
  specifying the lanes on which $L$ is to act, the
  \emph{XOR-bridge micro-kernel} executes:
  \begin{enumerate}[nosep,label=(\arabic*)]
    \item $w \leftarrow \mathrm{reg} \wedge \neg m$
      (witness: the lanes outside $L$'s range);
    \item $\mathrm{reg} \leftarrow L(\mathrm{reg})$
      (apply $L$ everywhere, including to the
      witness lanes, which will be corrected);
    \item $\mathrm{reg} \leftarrow \mathrm{reg}
      \oplus L(w) \oplus w$ (cancel $L$'s effect on
      the witness lanes: by $\GF(2)$-linearity
      $L(\mathrm{reg}) \oplus L(w) = L(\mathrm{reg}
      \oplus w) = L(\mathrm{reg}|_m)$, and $\oplus w$
      XORs the original witness values back in).
  \end{enumerate}
\end{definition}

\begin{proposition}[XOR-bridge correctness]
  \label{prop:xor-bridge}
  \footnote{Class \textsf{AHR}. Warrant:
  $\GF(2)$-linearity of $L$ and the fact that
  characteristic 2 makes $x \oplus x = 0$.
  Qualifier: $L$ must be strictly $\GF(2)$-linear on
  its entire domain; semi-linear or affine atoms do
  not qualify. Finite-$\kappa$ valid.}
  Let $L$ be $\GF(2)$-linear, $m$ a mask. After the
  three steps of Definition \ref{def:xor-bridge}, the
  register equals $L(\mathrm{reg}|_m) \oplus
  \mathrm{reg}|_{\neg m}$.
\end{proposition}

\begin{proof}
  Let $r_0$ be the register before step (1). Step (1)
  sets $w = r_0 \wedge \neg m = r_0|_{\neg m}$. Step
  (2) sets $\mathrm{reg} = L(r_0)$. Step (3):
  \begin{align*}
    \mathrm{reg}
    &= L(r_0) \oplus L(w) \oplus w \\
    &= L(r_0 \oplus w) \oplus w
      && \text{by $\GF(2)$-linearity of $L$} \\
    &= L(r_0|_m) \oplus r_0|_{\neg m}
  \end{align*}
  which is exactly what was claimed: $L$ applied on
  $m$, identity on $\neg m$.
\end{proof}

\begin{remark}[Free for linearity-preserving masks]
  When $L$ is the identity on $\neg m$ (i.e.\ $L$
  affects only $m$-lanes in the first place, as is
  the case for every atom with a strict input tile),
  step (3) collapses to $\mathrm{reg} \leftarrow
  \mathrm{reg} \oplus w \oplus w = \mathrm{reg}$, so
  the bridge is trivially a single XOR at the
  boundary, not three operations. The cost model
  credits this case via a production subclass
  \texttt{time\_share\_register\_disjoint}.
\end{remark}

\subsection{The \texttt{time\_share\_register} production}

\begin{definition}[\texttt{time\_share\_register}]
  \label{def:time-share}
  The \emph{time-share-register} production
  $\texttt{time\_share\_register}[L_1, L_2, \ldots,
  L_n]$ executes atoms $L_i$ sequentially on one
  shared register, with an XOR-bridge
  (Definition \ref{def:xor-bridge}) between each
  consecutive pair. The attribute combination is
  \begin{align}
    \VRAM &= \max_i \VRAM(L_i) + \VRAM_{\mathrm{I/O}} \\
    \SMEM &= \max_i \SMEM(L_i) + \SMEM_{\mathrm{masks}} \\
    \REG  &= \max_i \REG(L_i)  + 1 \quad \text{(witness register)} \\
    \LAT  &= \sum_i \LAT(L_i) + 2(n-1) \cdot \LAT_\oplus
  \end{align}
  where $\LAT_\oplus$ is the latency of one XOR
  operation.
\end{definition}

The cost combinator maxes $\REG$ rather than summing
it: at any given time only one $L_i$'s live set is in
the register, plus a single-word witness. This is the
exact opposite of the $\REG$-summing behavior of
\texttt{fused\_matrix\_apply}; the planner chooses
between them based on which resource is binding.

\subsection{Commutativity under lane-disjoint masks}

\begin{proposition}[Time-share commutativity]
  \label{prop:ts-commute}
  \footnote{Class \textsf{AHR}. Warrant: lane-disjoint
  masks make each $L_i$ a linear map on its own
  sub-register, independent of the others.
  Qualifier: atoms with overlapping masks do not
  commute. Finite-$\kappa$ valid.}
  If the atoms $L_1, \ldots, L_n$ of a
  \texttt{time\_share\_register} production have
  pairwise disjoint masks (Proposition
  \ref{prop:mask-disjoint}), then any permutation of
  the atoms produces the same final register state.
\end{proposition}

\begin{proof}
  Each $L_i$ acts on a sub-register determined by
  $m_i$. Since $m_i \wedge m_j = 0$, the action of
  $L_i$ does not modify the domain of $L_j$. Hence
  $L_1 \circ L_2 = L_2 \circ L_1$ as operations on
  the packed register, and by induction on permutation
  length, any permutation yields the same result.
\end{proof}

\begin{remark}
  Commutativity licenses the parser to collapse
  permutation-equivalent time-share productions
  during chart construction, mirroring
  Proposition \ref{prop:xor-commute}. For $n$
  lane-disjoint atoms, this shrinks the chart by a
  factor of $n!$.
\end{remark}

\subsection{Non-disjoint masks}

When masks overlap, the XOR-bridge still applies
per-atom but commutativity is lost: successive atoms
may act on overlapping lanes, and the order of their
applications affects the final state.

\begin{proposition}[Ordered time-share]\label{prop:ts-ordered}
  For overlapping masks, the
  \texttt{time\_share\_register} production is
  ordered: the plan's semantics agrees with the
  naive sequential execution of the atoms in the
  stated order, using Proposition \ref{prop:xor-bridge}
  at each step.
\end{proposition}

\begin{proof}
  By induction on the number of atoms. Each XOR-bridge
  step leaves the non-masked lanes untouched by
  Proposition \ref{prop:xor-bridge}, so the next
  atom sees the same register state it would in a
  naive execution.
\end{proof}

\subsection{Non-linear atoms break the bridge}

\begin{remark}[Non-linearity as a wall]
  The XOR-bridge requires $L$ to be $\GF(2)$-linear.
  Atoms like $\mathrm{pcnt}_2$, which reduce an
  $n$-bit input to a single bit, are not linear in
  the required sense. They cannot participate in a
  \texttt{time\_share\_register} production; instead,
  the planner inserts a \texttt{spill\_register}
  boundary (Section \ref{sec:parsing}) that
  materializes the register to shared memory before
  the non-linear atom and re-loads afterwards. The
  cost of spill is explicit in the attribute rule
  and is subject to the usual budget enforcement.
\end{remark}

\begin{conjecture}[Semi-reversibility]
  \label{conj:nonlinear}
  Certain nominally non-linear atoms ($\mathrm{pcnt}_2$,
  comparison, reduce-sum-mod-$p$) admit a
  \emph{paired inverse atom} that undoes the
  information loss in the context of a specific
  surrounding workload. A formal characterization
  of when such pairs exist, and how they compose
  under an extended time-share production, is open
  (\S\ref{sec:reading-guide}, X2).
\end{conjecture}

\subsection{Nested productions}

A \texttt{time\_share\_register} production may have
\texttt{fused\_matrix\_apply} nodes as its children, or
vice versa. The nesting is unconstrained: the cost
attribute rule for each parent is applied to its
immediate children's aggregate cost, regardless of
how those children were constructed.

\begin{proposition}[Nested cost additivity]
  \label{prop:nest-cost}
  For a production $P$ with children $C_1, \ldots,
  C_k$ whose cost vectors are $\bar{c}_i$ (each itself
  the product of nested productions), the parent's
  cost under the stated combinator is
  $\mathrm{combinator}_P(\bar{c}_1, \ldots,
  \bar{c}_k) + \bar{c}_{P_0}$, where $\bar{c}_{P_0}$
  is the production's baseline.
\end{proposition}

This is a direct consequence of the attribute grammar
being well-defined compositionally, but we state it
explicitly because it is the property that makes the
chart parser's dynamic programming valid: each chart
item's cost is computable from its children's costs
in one step.

\subsection{Putting the two compound productions together}

Section \ref{sec:fusion} introduced spatial fusion
(atoms concurrent, $\REG$ summed); this section has
introduced temporal multiplexing (atoms sequential,
$\REG$ maxed). The planner has both in its toolkit
and will choose between them (or nest them) based on
which resource in the budget vector $B$ is binding.
Under $\LAT$-binding budgets, \texttt{fused\_matrix\_apply}
dominates; under $\REG$-binding,
\texttt{time\_share\_register} dominates; the middle
ground of moderate $\LAT$ and $\REG$ pressure is
typically won by nested plans where several small
fused blocks are chained via XOR-bridges into a
single sequential kernel. The mechanics of that
choice is the subject of Section \ref{sec:parsing}.
